% !PRINTED 637
\weberpage{637}
%<<CONT
wo $m = n'n''$ ist, und die $x_{\nu,\nu'}$ ganze rationale Funktionen von $z$
sind. Aus (7) ergibt sich durch Elimination der $\omega_i$:

\[
\begin{vmatrix}
x_{1,1} - \omega, & x_{1,2}, & \ldots, & x_{1,m} \\
x_{2,1}, & x_{2,2} - \omega, & \ldots, & x_{2,m} \\
\hdotsfor{4} \\
x_{m,1}, & x_{m,2}, & \ldots, & x_{m,m} - \omega
\end{vmatrix} = 0,
\]

und dies ist eine Gleichung für $\omega$ von der Form (5).

Durch wiederholte Anwendung dieses Satzes folgt, daß jede
\gesperrt{ganze rationale Funktion von ganzen Funktionen wieder
eine ganze Funktion ist.}

7. \gesperrt{Eine ganze Funktion $\omega$ heißt durch eine ganze
Funktion $\omega'$ teilbar, wenn eine dritte ganze Funktion
$\omega''$ existiert, so daß}

\weq{8}{\omega = \omega'\omega''}

\gesperrt{ist.}

Aus dieser Definition ergibt sich ohne weiteres:

\gesperrt{Ist $\omega$ teilbar durch $\omega'$ und $\omega'$ teilbar durch $\omega''$, so
ist auch $\omega$ durch $\omega''$ teilbar.}

\gesperrt{Sind $\omega'$ und $\omega''$ durch $\omega$ teilbar, so ist auch $\omega' \pm \omega''$
durch $\omega$ teilbar.}

\gesperrt{Sind $\omega_1, \omega_2, \omega_3, \ldots$ durch $\omega$ teilbar und $\omega_1', \omega_2', \omega_3', \ldots$
beliebige ganze Funktionen, so ist auch}

\[
\omega_1 \omega_1' + \omega_2 \omega_2' + \omega_3 \omega_3' + \cdots
\]

\gesperrt{durch $\omega$ teilbar.}

\begin{center}
\textbf{§ 176. Minimalbasis und Körperdiskriminante.}
\end{center}

Da man jede Funktion des Körpers $\Omega$ durch Multiplikation
mit einer ganzen rationalen Funktion in eine ganze algebraische
Funktion von $z$ verwandeln kann, so gibt es auch Körperbasen,
die aus \gesperrt{ganzen Funktionen} bestehen (z. B. nach der Bezeichnung
 in § 171, (2) die Potenzen von $a_0 \theta$).

Ist nun

\weq{1}{\omega_1, \ \omega_2, \ \ldots, \ \omega_n}

eine solche aus ganzen Funktionen bestehende Basis, so ist jede
in der Form

\weq{2}{\omega = x_1 \omega_1 + x_2 \omega_2 + \cdots + x_n \omega_n}
%CONT>>
